\scanpage{72}
On pose
\[
\left\{\begin{aligned}
\mathrm{cn}_1^2z&=1-(e_2-e_1)\mathrm{sn}_1^2z=\frac{\unsure{\wp|p}z-e_2}{\wp z-e_1}&(\mathrm{cn}_1\,0=1)\\
\mathrm{dn}_1^2z&=1-(e_3-e_1)\mathrm{sn}_1^2z=\frac{\wp z-e_3}{\wp z-e_1}&(\mathrm{dn}_1\,0=1)
\end{aligned}\right.\tag*{5}
\]
\[
\left\{\begin{aligned}
\mathrm{cn}_1z&=\frac{\sigma\omega_1}{\sigma\omega_2}e^{(\eta_2-\eta_1)z}\frac{\sigma(z-\omega_2)}{\sigma(z-\omega_1)}
=\frac{\sigma\omega_1}{\sigma\omega_2}e^{-(\eta_2-\eta_1)z}\frac{\sigma(z+\omega_2)}{\sigma(z+\omega_1)}\\
\mathrm{dn}_1z&=\frac{\sigma\omega_1}{\sigma\omega_3}e^{(\eta_3-\eta_1)z}\frac{\sigma(z-\omega_3)}{\sigma(z-\omega_1)}
=\frac{\sigma\omega_1}{\sigma\omega_3}e^{-(\eta_3-\eta_1)z}\frac{\sigma(z+\omega_3)}{\sigma(z+\omega_1)}
\end{aligned}\right.\tag*{6}
\]

\deleted{\uline{Poles de $\mathrm{cn}_1z$}}
\[
\left\{\begin{aligned}
\mathrm{cn}_1\,{-z}&=\mathrm{cn}\,z &\mathrm{dn}_1\,{-z}&=\mathrm{dn}_1z\\
\mathrm{cn}_1(z+2\omega_1)&=-\mathrm{cn}\,z &\mathrm{dn}_1(z+2\omega_1)&=-\mathrm{dn}_1z\\
\mathrm{cn}_1(z+2\omega_2)&=-\mathrm{cn}\,z &\mathrm{dn}_1(z+2\omega_{3})&=-\mathrm{dn}_1z\\
\Vert\ \mathrm{cn}_1(z+2\omega_3)&=\mathrm{cn}\,z &\mathrm{dn}_1(z+2\omega_2)&=\mathrm{dn}_1z\\
\mathrm{cn}_1(z+\deleted{2}4\omega_1)&=\mathrm{cn}\,z &\mathrm{dn}_1(z+4\omega_3)&=\mathrm{dn}_1z\\
\mathrm{cn}_1(z+4\omega_2)&=\mathrm{cn}\,z &\mathrm{dn}_1(z-4\omega_1)&=\mathrm{dn}_1z
\end{aligned}\right.\tag*{7}
\]
$\mathrm{cn}_1z$ et $\mathrm{dn}_1z$ sont paires

$\mathrm{cn}_1z$ et $\mathrm{dn}_1z$ ont resp\textsuperscript{t} les périodes $2\omega_3,4\omega_1,4\omega_3$ et $2\omega_2,4\omega_3,4\omega_{\unsure{1|2}}$ ; elles sont d'ordre deux.

\uline{Poles de $\mathrm{cn}_1z$} : homologues de $\omega_1$, \deleted{résidu $\unsure{z\omega_1+2\omega_3|\omega_1+2\omega_3|3\omega_1+2\omega_3}$} \inserted{et} $-\omega_1$

Si la disposition des vecteurs $\omega_1,\omega_2$ est directe, le résidu en $\omega_1$ est
\[
\text{résidu en }\omega_1:\ \frac{i}{k_1\mathrm{sn}_1\omega_2}\ ;=\pm i\sqrt{e_3-e_1}
\]
\uline{Zéros de $\mathrm{cn}_1z$} homologues de $+\omega_2$ et $-\omega_2\deleted{+\unsure{2\omega_1}}$ \deleted{\unsure{ou $\omega_2+2\omega_3$|ou $\omega_1+2\omega_3$}} ;

\uline{Poles de $\mathrm{dn}_1z$} : homologues de $\omega_1$ et $-\deleted{\unsure{3|2}}\omega_1$ ; avec l'hypothèse précédente, on a
\[
\text{résidu en }\omega_1:\ \frac{-i}{k_1\mathrm{sn}_1\omega_3}=\pm i\sqrt{e_2-e_1}
\]
\uline{Zéros de $\mathrm{dn}_1z$} : homologues de $\omega_3$ et $-\omega_3\deleted{+\unsure{2\omega_2}}$

% archive sheet number: 71
